\documentclass{beamer}
\usepackage{amsmath}
\setbeameroption{show notes on second screen=right}
\setbeamertemplate{note page}[plain]

\begin{document}

\begin{frame}
\frametitle{}

When multiplying fractions, there is no requirement for a common denominator.
\note[item]{\alert<1>{When multiplying fractions, there is no requirement for a common denominator.}}
\note[item]{\alert<2>{To multiply fractions, we just multiply straight across.}}

\vskip10pt

\pause 
\pause 
{\bf Example.} $\frac25 \times \frac37$.
\note[item]{\alert<3>{Let's look at an example. What is two fifths times three sevenths?}}

\vskip4pt

\pause $=\frac{2 \cdot 3}{5 \cdot 7}$ 
\note[item]{\alert<4>{To multiply fractions, we multiply the numerators together and the denominators together.}}

\vskip4pt

\pause $=\frac{6}{35}$
\note[item]{\alert<5>{Simplifying the top and bottom, our single fraction is six over thirty-five.}}

\end{frame}

\begin{frame}
\frametitle{}

{\bf Example.} $\frac{1}{10} \times \frac{1}{6}$. 
\note[item]{\alert<1>{Let's look at another example. What is one tenth times one sixth?}}

\vskip10pt

\pause 
\note[item]{\alert<2>{As a reminder, for multiplying fractions, there is no need for a common denominator.}}

\vskip4pt

\pause $=\frac{1 \cdot 1}{10 \cdot 6}$
\note[item]{\alert<3>{To multiply these fractions, we multiply the numerators together and the denominators together.}}

\vskip4pt

\pause $=\frac{1}{60}$
\note[item]{\alert<4>{Simplifying the top and bottom, our single fraction is one over sixty.}}


\end{frame}


\begin{frame}
\frametitle{}

{\bf Example.} $\frac{3}{7} \times \frac{2}{7}$. 
\note[item]{\alert<1>{Let's look at another example. What is three sevenths times two sevenths?}}

\vskip10pt

\pause
\note[item]{\alert<2>{Notice both fractions have seven as a denominator. That should not change our process. We will still multiply straight across.}}

\pause $=\frac{3 \cdot 2}{7 \cdot 7}$ \pause
\note[item]{\alert<3>{We multiply numerators three and two together.}}
\note[item]{\alert<4>{We also multiply the denominators seven and seven together.}}

\vskip4pt
\pause $=\frac{6}{49}$
\note[item]{\alert<5>{Simplifying the numerator and denominator, we have six over forty-nine.}}
\note[item]{\alert<6>{The incorrect belief that you need a common denominator to multiply fractions leads people to get the incorrect final answer of six over seven.}}
\note[item]{\alert<7>{However, we only need a common denominator for adding and subtracting fractions, not multiplying.}}
\note[item]{\alert<8>{The correct answer is six over forty-nine. The denominator of forty-nine is the result of multiplying the first denominator (which is seven) by the second denominator (which is also coincidentally seven).}}

\end{frame}


%% Blank frames at the end to prevent weird shifts.
\begin{frame}
\end{frame}
\begin{frame}
\end{frame}

\end{document}
